\textbf{Platforms.} Table~\ref{tab:platforms} lists the two machines. The \emph{Spartan GPU node} is the reference platform: it has both a multicore CPU and a GPU, so serial, OpenMP, CUDA and hybrid runs are compared on the same machine, one benchmark job per node under Slurm~\cite{yoo2003slurm}. The \emph{Apple M4} laptop has no NVIDIA GPU; it is used only for the CPU implementations and is never mixed with GPU-node numbers. The M4 has 4 performance and 6 efficiency cores, so its thread-scaling curve is not that of a homogeneous machine.

\begin{widetable}{The two machines. ``Physical cores / logical threads'' are the hardware counts; ``OpenMP threads used'' is what the benchmark could use. The Spartan job was allocated 24 of the node's 28 CPUs.}{tab:platforms}
\begin{tabular*}{\textwidth}{@{\extracolsep{\fill}}p{0.17\textwidth}p{0.36\textwidth}p{0.39\textwidth}@{}}
\toprule
 & \textbf{Spartan GPU node (reference)} & \textbf{Apple M4 laptop (CPU only)} \\
\midrule
CPU & 2$\times$ Intel Xeon E5-2680 v4, 2.40\,GHz & Apple M4 \\
Physical cores / logical threads & 28 / 28 (14 per socket, 1 thread per core) & 10 / 10 (4 performance + 6 efficiency) \\
OpenMP threads used & 1, 2, 4, 8, 12, 16, 24 & 1, 2, 4, 6, 8, 10 \\
Cache & 32\,KiB L1d and 256\,KiB L2 per core; 35\,MiB L3 per socket; 2 NUMA nodes & 16\,MiB L2 (performance cluster), 4\,MiB L2 (efficiency cluster) \\
Memory & 125\,GiB & 16\,GiB \\
GPU & Tesla K40m: compute capability 3.5, 15 SMs, 11{,}441\,MiB, 384-bit bus, memory clock 3{,}004\,MHz & none \\
Operating system & Linux 5.14 (Slurm, partition \texttt{gpus}) & macOS 26.6.2 \\
CUDA & \texttt{nvcc} 11.7.99, runtime 11.7, driver 470.256.02 (CUDA 11.4); \texttt{-arch=sm\_35} & not available \\
Host compiler & g++ 11.4.1, OpenMP 4.5; \texttt{-O3 -std=c++14}, plus \texttt{-march=native} for the reference runs & Apple clang 21.0.0 with libomp 23.1.2; \texttt{-O3 -std=c++14} \\
\bottomrule
\end{tabular*}
\end{widetable}


\textbf{Protocol.} For every configuration (implementation, batch size, thread count, block size, number of streams) we run 3 unrecorded warm-up passes and then \val{Sp:runs} recorded passes, and store every pass as one row of a CSV file together with the stage times, the accuracy of that pass, and the hardware and compiler description~\cite{hoefler2015scientific}. Raw files are never overwritten. A pass processes the 10{,}000 test images; the timers start after all host buffers have been allocated, and one-time set-up (device buffers and the weight upload) is excluded for every implementation. OpenMP threads are bound with \texttt{OMP\_PROC\_BIND=close} and \texttt{OMP\_PLACES=cores} on Linux; macOS ignores thread binding.

\textbf{Timing.} \emph{End-to-end} time is the wall clock of a whole pass. CPU stages are timed with the steady clock; the GPU stages (host-to-device copy, the three kernels, device-to-host copy) are timed with CUDA events~\cite{harris2012metrics}. Kernel-only time is always reported together with end-to-end time, because they answer different questions: the first measures how fast the GPU computes, the second what an application actually waits for~\cite{gregg2011data}. In the pipelined hybrid the GPU stages overlap each other and the CPU work, so their sum can exceed the end-to-end time.

\textbf{Metrics.} With $T$ the end-to-end time of a pass: speedup $S=T_{\mathrm{serial}}/T_{\mathrm{impl}}$ on the same machine and batch size; OpenMP efficiency $E=S/p$ for $p$ threads; throughput $=10{,}000/T$ images per second; transfer share $=(T_{\mathrm{H2D}}+T_{\mathrm{D2H}})/T$. Each mean is given with a 95\,\% Student-$t$ confidence interval. Hybrid variants are compared with Welch's $t$-test and a bootstrap interval of the ratio of means (5{,}000 resamples). The Amdahl serial fraction $f$ is fitted by least squares to the measured speedups of each batch size, and the Karp--Flatt metric is computed for every thread count. When we say ``best'' thread count we mean the tested count with the lowest mean time.

\textbf{Compiler flags.} All CPU code is compiled with \texttt{-O3}. On the Spartan node, whose default x86-64 target only uses 128-bit vector instructions, the reference runs use \texttt{-march=native} so that the compiler may use the vector instructions of the host CPU for these straightforward loops~\cite{lee2010debunking}; a second set of runs with plain \texttt{-O3} quantifies how much this matters. No implementation uses a tuned BLAS routine or a hand-tiled kernel.
